\section{Monica's Farewell and Augustine's Tears}

In the ninth book of his \emph{Confessions}, Augustine remembers a conversation with his mother shortly before her death. They stand together at a window in Ostia and speak of the eternal life of the saints (IX.10.23--26). Their minds ascend through created things toward the Wisdom by whom all things are made. Then they return to ordinary speech. The account gives a moment of shared contemplation within a particular mother and son's history, after the long years in which Monica had prayed for his conversion.

Monica tells him that her great earthly hope has been fulfilled: she has seen him become a Catholic Christian and a servant of God. Soon afterward she falls ill. When Augustine's brother expresses the wish that she might die in her own country, she sets aside anxiety about the place of burial. Her request is to be remembered at the Lord's altar. Augustine marvels at the change, because she had formerly cared deeply about being buried beside her husband (IX.11.27--28). Grace has loosened a cherished attachment without diminishing the love beneath it.

Monica trusts that God will know where to raise her body. She can therefore accept burial far from home and ask for what she most desires: remembrance at the altar. Her changed wish shows how faith in resurrection can loosen anxiety about circumstances we cannot control. We can make appropriate arrangements and honor family bonds, yet still entrust to God what illness, distance, or a sudden death prevents us from arranging. Those who survive us can pray for us wherever we are buried, as Monica asks her sons to do.

Augustine's account of the death itself is painfully exact. He closes her eyes. Grief rises within him, and he attempts to restrain it. Those who hear him speak with others may suppose he is not sorrowful, while he knows the inward wound of their life together having been torn apart (IX.12.29--32). His theological understanding has not prevented the loss from reaching the habits of love. The familiar presence, built from years of shared life, has been taken away.

He also finds himself sorrowing over his sorrow. He is troubled that ordinary human feeling has such power over him. Later, after rest, he recalls Monica's tenderness and service and allows himself to weep before God (IX.12.33). The tears he had restrained become a place of relief. He does not invite the reader to mock them. He asks for the charity that recognizes a fellow human being grieving for the mother who had once wept for his salvation.

St.~Ambrose likewise begins the second book of \emph{On the Death of His Brother Satyrus} by remembering the fresh wound to which the first address had given expression. He says that remedies applied too sharply may increase pain, and allows a place to natural affection and tears before turning toward a wider consolation in resurrection (II.1--3). Later he asks that memory become grateful rather than something we must avoid because it only hurts (II.8). The bishop speaks as a brother. His exhortation toward hope comes from within the bond whose severance he feels.

The apostolic command is that we not grieve as those without hope (1~Thess 4:12--17 in the Douay--Rheims). Hope changes the meaning and destination of grief. It does not require the bereaved to demonstrate, at once, that every ache has subsided. Tears can accompany confidence in Christ. Remembering a gift can remain painful while gratitude slowly grows. The Church gives us words and actions in which both sorrow and hope can be carried when our private strength is insufficient.

Augustine also prays for Monica's sins. He praises the grace visible in her life, yet acknowledges her need for forgiveness after Baptism (IX.13.34--36). Her good works are God's gifts; she depends on the mercy of the Christ who redeemed her. He trusts the innocent blood shed by the Savior and asks that his mother remain under that protection. The son who remembers her holiness still intercedes for her, relying on the same mercy she had sought for him.

At the close of the book he asks its readers to remember Monica and Patricius at the altar (IX.13.37). He seeks the prayers of people who will know his parents through his account of them. We can follow him in asking others to pray for our own dead. Affection alone cannot tell us their eternal state. It gives us reason to desire their salvation and to continue asking mercy for them, through the intercession the Church teaches and especially at the altar where Monica wished to be remembered.

At a Christian funeral the Church honors the body, proclaims resurrection, and commends the deceased to God. When the Eucharistic sacrifice is celebrated, she asks the Father to purify the departed and admit them to the fullness of his kingdom (CCC 1684--1690). The funeral therefore helps the dead as well as consoling the living. Our desire to express admiration and remember the person's gifts should leave room for this intercession. Even when a life has given much reason for gratitude, the person still needs the mercy sought in the Church's prayer.

This prayer continues after the funeral. The Church's suffrages, especially the Eucharistic sacrifice, help those undergoing the final purification of the saved (CCC 1030--1032). We can remember the dead at Mass, pray for them, and perform works of charity in their memory. A visit to a grave may renew sorrow, but it can also become an occasion to commend the person to God again. As time passes, anniversaries and familiar prayers help us continue an intercession we might otherwise neglect when the first intensity of bereavement has subsided.

